CLASE 2

formato de diapositivas como el ejemplo (toda la documentacion es word y powerpoint sin cambiar el formato)

y todos los documentos en formato word (ver ejemplo)

autores:
apellidos nombres (oredn alf)


bajo nivel:
patrones de disenio 




alto nivel:
patrones arqutectonicos


patrones arquitectonicos orientados a la vista e interfaz
mvc (todo en mvc para esta clase)
mmvc

ejemplo:
login -> en el servidor

Lo promero que hce el ingeniero o arquitecto de software es la planeacion y esta comienza con un ERS:

Documento ERS se hace con la IEES 830 con una variacion del ingeniero moster
RF -> que se hara; del lado del negocio
RNF -> como se lo hara involucra todo el estack tecnologico y las razon del porque de las mismas, asi mismo las metodologias de desarooolo;
   del lado delas tecnologias de la informacion. 

La arquitectura es darle calidad al codigo para que sea mas barato mantener y facil de escalar.

El primer patron de disenio deberia ser el singleton para la conexion (pool de conexciones) ya que es una solo ainstancia 
de esta conexion. 

api graphql vs soap vs rest
Característica	REST	GraphQL	SOAP
Tipo	Estilo arquitectónico	Lenguaje de consulta para APIs	Protocolo
Formato típico	JSON	JSON	XML
Endpoints	Varios	Normalmente uno	Normalmente uno
Métodos	GET, POST, PUT, PATCH, DELETE	Principalmente POST	Principalmente POST
Contrato	OpenAPI opcional	Schema obligatorio	WSDL
Flexibilidad de respuesta	Media	Muy alta	Baja
Facilidad	Alta	Media	Baja
Overfetching	Puede ocurrir	Lo evita	Puede ocurrir
Seguridad empresarial	Buena	Buena	Muy completa / WS-Security
Uso actual típico	APIs web y microservicios	Frontends con datos complejos	Sistemas empresariales/legacy

para aplicar arquitectura se abtrae todas las funciones debajo de una interfaz
todo lo que se guarda se abstrae mediante un repositorio
(usar estructuras mas abstractas, separadas por responsabilidades)

Jakarta Enterprise Edition (Jakarta EE) es una plataforma para desarrollar aplicaciones empresariales en Java.
.NET Core

(todas las capas, que hacen, etc de memoria)

reviddsar (material de revision *java y .NET  capas de puntonet como funcionan)